\begin{examplebox}
	\textbf{\textcolor{CExample}{例1（基础）}}

	\textbf{\textcolor{CExample}{题目：}}
	已知函数$f(x)$满足$f(x+6)=f(x-6)$对任意$x\in\mathbb{R}$成立，求$f(x)$的周期。

	\textbf{\textcolor{CExample}{解题步骤：}}
	\begin{description}[style=nextline, leftmargin=2.8em, labelsep=0.5em, itemsep=0.45em]
		\item[\textbf{第一步（识别条件）：}] 条件形如$f(x+a)=f(x-a)$，这里$a=6$。
			\par\textbf{理由：} 直接对比已知形式。
		\item[\textbf{第二步（应用结论）：}] 由周期结论，周期$T=2a=12$。
			\par\textbf{理由：} 若$f(x+a)=f(x-a)$恒成立，则$2a$是周期。
		\item[\textbf{第三步（下结论）：}] 因此$f(x)$的一个周期为$12$。
			\par\textbf{理由：} 满足周期定义$f(x+12)=f(x)$。
	\end{description}

	\textbf{\textcolor{CExample}{关键结论：}}
	\textbf{$T=12$。}

	\medskip
	\textbf{\textcolor{CExample}{例2（稍变形）}}

	\textbf{\textcolor{CExample}{题目：}}
	已知$f(x)$满足$f(x+2)=f(x-2)$，且$f(1)=3$，求$f(9)$的值。

	\textbf{\textcolor{CExample}{解题步骤：}}
	\begin{description}[style=nextline, leftmargin=2.8em, labelsep=0.5em, itemsep=0.45em]
		\item[\textbf{第一步（求周期）：}] 条件中$a=2$，所以周期$T=2a=4$。
			\par\textbf{理由：} 周期判定结论直接得出。
		\item[\textbf{第二步（周期转化）：}] 因为$9=1+2\times4$，所以$f(9)=f(1+2\times4)=f(1)=3$。
			\par\textbf{理由：} 周期函数有$f(x+4k)=f(x)$，$k\in\mathbb{Z}$。
		\item[\textbf{第三步（得出答案）：}] 故$f(9)=3$。
			\par\textbf{理由：} 代入已知$f(1)=3$。
	\end{description}

	\textbf{\textcolor{CExample}{关键结论：}}
	\textbf{$f(9)=3$。}

	\medskip
	\textbf{\textcolor{CExample}{例3（综合应用）}}

	\textbf{\textcolor{CExample}{题目：}}
	定义在$\mathbb{R}$上的函数$f(x)$满足$f(x+3)=f(x-3)$，且当$x\in[-3,3)$时$f(x)=x^2$，求$f(2024)$的值。

	\textbf{\textcolor{CExample}{解题步骤：}}
	\begin{description}[style=nextline, leftmargin=2.8em, labelsep=0.5em, itemsep=0.45em]
		\item[\textbf{第一步（求周期）：}] 条件中$a=3$，所以周期$T=6$。
			\par\textbf{理由：} 周期结论。
		\item[\textbf{第二步（转化自变量）：}] 计算$2024\div6$余数：$2024=337\times6+2$，故$f(2024)=f(2)$。
			\par\textbf{理由：} $f(x+6k)=f(x)$，$k\in\mathbb{Z}$。
		\item[\textbf{第三步（代入解析式）：}] 因为$2\in[-3,3)$，所以$f(2)=2^2=4$。
			\par\textbf{理由：} 已知在$[-3,3)$上$f(x)=x^2$。
	\end{description}

	\textbf{\textcolor{CExample}{关键结论：}}
	\textbf{$f(2024)=4$。}
\end{examplebox}
